% --- INSTITUTIONAL TEMPLATE FED v7.4.9 ---
% Estándar de Reporteo APA 7 (Babel-Free)
\documentclass[12pt,spanish]{article}

\usepackage{graphicx}
\usepackage{geometry}
\usepackage{indentfirst} 
\usepackage{caption}
\usepackage{floatrow} 

% Forzar estilos de caption antes de cualquier carga
\floatsetup[figure]{capposition=top}
\floatsetup[table]{capposition=top}

\usepackage{booktabs}
\usepackage{float}
\usepackage{longtable}
\usepackage{threeparttable}
\usepackage{array}
\usepackage{calc}
\usepackage{amsmath, amssymb}
\usepackage{fancyhdr}
\usepackage{mathptmx} 
\usepackage[T1]{fontenc}
\usepackage{setspace} 
\usepackage{ragged2e} 
\usepackage{titlesec}
\usepackage{url}
\def\UrlBreaks{\do\/\do-\do\.\do\_\do\?\do\=\do\&\do\%}
\usepackage[hidelinks,breaklinks]{hyperref}

% --- Pandoc Compatibility Block (Iron Clad) ---
\providecommand{\tightlist}{%
  \setlength{\itemsep}{0pt}\setlength{\parskip}{0pt}}

\usepackage{xcolor}
\usepackage{fancyvrb}
\usepackage{framed}

% --- Code Highlighting (Pandoc Standard) ---
\AtBeginDocument{
  \definecolor{shadecolor}{RGB}{248,248,248}
}

% Compatibilidad Pandoc 3.x
\ifdefined\pandocbounded\else
  \newcommand{\pandocbounded}[1]{#1}
\fi

% --- Pandoc Citeproc (CSL) compatibility ---
\newlength{\cslhangindent}
\setlength{\cslhangindent}{1.27cm} 
\newlength{\csllabelwidth}
\setlength{\csllabelwidth}{3em}
\newcommand{\citeproctext}{} 
\newcommand{\citeprocitem}[2]{\item} 
\newenvironment{CSLReferences}[2] 
 {\begin{list}{}{
  \setlength{\leftmargin}{\cslhangindent}
  \setlength{\parsep}{0pt}
  \setlength{\itemsep}{6pt}
  \ifodd #1
   \setlength{\leftmargin}{\leftmargin + \csllabelwidth}
   \setlength{\itemindent}{-\csllabelwidth}
  \fi
  }
  \renewcommand{\bibitem}[2][]{\item} % Elimina los corchetes [ ]
 }
 {\end{list}}
\newcommand{\CSLBlock}[1]{#1\hfill\break}
\newcommand{\CSLLeftMargin}[1]{\parbox[t]{\csllabelwidth}{#1}}
\newcommand{\CSLRightInline}[1]{\parbox[t]{\linewidth - \csllabelwidth}{#1}\break}
\newcommand{\CSLIndent}[1]{\hspace{\cslhangindent}#1}

\makeatletter
\@ifundefined{paragraph}{\newcounter{paragraph}}{}
\@ifundefined{subparagraph}{\newcounter{subparagraph}}{}
\makeatother

% Solución al error 'No counter none defined' en longtable
\makeatletter
\newcounter{none}
\@ifpackageloaded{caption}{
  \DeclareCaptionType{none}
  \captionsetup[none]{name=Tabla}
}{}
\makeatother

\hypersetup{
  colorlinks=true,
  linkcolor=blue,
  filecolor=blue,
  citecolor=blue,
  urlcolor=blue
}

% --- Macros Institucionales ---
\newcommand{\fuente}[1]{
  \vspace{4pt}
  \begin{flushleft}
    \small
    \textit{Nota.} #1
  \end{flushleft}
}

% --- Localización Manual Determinista (Sin Babel) ---
\AtBeginDocument{
  \renewcommand{\tablename}{Tabla}
  \renewcommand{\figurename}{Figura}
  \renewcommand{\contentsname}{Índice}
  \renewcommand{\listtablename}{Índice de Tablas}
  \renewcommand{\listfigurename}{Índice de Figuras}
  \renewcommand{\abstractname}{Resumen}
  \renewcommand{\refname}{Referencias}
  
  \RaggedRight
  \setlength{\parindent}{1.27cm}
  \setlength{\parskip}{0pt}
}

\DeclareCaptionLabelSeparator{newline}{\\ \vspace{2pt}}
\captionsetup{
  position=top,
  justification=RaggedRight,
  singlelinecheck=false,
  labelfont=bf,
  textfont=it,
  labelsep=newline,
  font=small,
  skip=10pt
}
\captionsetup[figure]{name=Figura}
\captionsetup[table]{name=Tabla}

% Márgenes APA 7
\geometry{
  a4paper,
  margin=25.4mm,
  headheight=15mm
}

% --- Encabezados ---
\pagestyle{fancy}
\fancyhf{} 
\fancyhead[L]{\includegraphics[height=1.2cm]{\logoUNL}}
\fancyhead[R]{\includegraphics[height=1.2cm]{\logoECONOMIA}}
\renewcommand{\headrulewidth}{0pt}
\fancyfoot[R]{\thepage}

\fancypagestyle{plain}{
  \fancyhf{}
  \fancyhead[L]{\includegraphics[height=1.2cm]{\logoUNL}}
  \fancyhead[R]{\includegraphics[height=1.2cm]{\logoECONOMIA}}
  \fancyfoot[R]{\thepage}
  \renewcommand{\headrulewidth}{0pt}
}

% --- Títulos ---
\titleformat{\section}{\normalfont\normalsize\bfseries\centering}{\thesection}{1em}{}
\titleformat{\subsection}{\normalfont\normalsize\bfseries}{\thesubsection}{1em}{}
\titleformat{\subsubsection}{\normalfont\normalsize\bfseries\itshape}{\thesubsubsection}{1em}{}
\titleformat{\paragraph}{\normalfont\normalsize\bfseries}{\theparagraph}{1em}{}
\titlespacing*{\section}{0pt}{12pt}{6pt}
\titlespacing*{\subsection}{0pt}{12pt}{6pt}
\titlespacing*{\subsubsection}{0pt}{12pt}{6pt}
\titlespacing*{\paragraph}{0pt}{12pt}{6pt}

\newenvironment{referenciasAPA}{
  \setlength{\parindent}{-1.27cm}
  \setlength{\leftskip}{1.27cm}
  \setlength{\parskip}{6pt}
}{\par}

\newcommand{\logoUNL}{/home/erick-fcs/Capital_Workstation/capital-workstation-libs/src/ecs_quantitative/reporting/assets/logos/logo_unl.png}
\newcommand{\logoECONOMIA}{/home/erick-fcs/Capital_Workstation/capital-workstation-libs/src/ecs_quantitative/reporting/assets/logos/logo_unl.png}\makeatletter\@ifpackageloaded{caption}{}{\usepackage{caption}}\AtBeginDocument{%
\ifdefined\contentsname
  \renewcommand*\contentsname{Table of contents}
\else
  \newcommand\contentsname{Table of contents}
\fi
\ifdefined\listfigurename
  \renewcommand*\listfigurename{List of Figures}
\else
  \newcommand\listfigurename{List of Figures}
\fi
\ifdefined\listtablename
  \renewcommand*\listtablename{List of Tables}
\else
  \newcommand\listtablename{List of Tables}
\fi
\ifdefined\figurename
  \renewcommand*\figurename{Figure}
\else
  \newcommand\figurename{Figure}
\fi
\ifdefined\tablename
  \renewcommand*\tablename{Table}
\else
  \newcommand\tablename{Table}
\fi
}\@ifpackageloaded{float}{}{\usepackage{float}}
\floatstyle{ruled}
\@ifundefined{c@chapter}{\newfloat{codelisting}{h}{lop}}{\newfloat{codelisting}{h}{lop}[chapter]}
\floatname{codelisting}{Listing}\newcommand*\listoflistings{\listof{codelisting}{List of Listings}}\makeatother\makeatletter\makeatother\makeatletter\@ifpackageloaded{caption}{}{\usepackage{caption}}
\@ifpackageloaded{subcaption}{}{\usepackage{subcaption}}\makeatother

\begin{document}

\begin{center}
    {\bfseries \Large UNIVERSIDAD NACIONAL DE LOJA}\\[0.3cm]
    {\bfseries \large FACULTAD JURÍDICA, SOCIAL Y
ADMINISTRATIVA}\\[0.2cm]
    {\bfseries CARRERA DE TRABAJO SOCIAL}\\[1.5cm]
    
    {\bfseries \large FUNDAMENTOS METODOLÓGICOS DEL TRABAJO
SOCIAL}\\[1cm]
    
    {\bfseries \Large Informe de Análisis sobre Corrientes Sociológicas
y Trabajo Social}\\[1.5cm]
\end{center}

\begin{center}
    \setstretch{1.1}
                        Erick Fabricio Condoy Seraquive \\
                \vspace{0.5cm}
    Lcda. Ligia Geomara Párraga Vélez \\ 
    
\end{center}

\vspace{1cm}

\doublespacing
\setlength{\parindent}{1.27cm}

\subsection{1. Introducción}\label{introducciuxf3n}

El Trabajo Social es una disciplina y profesión cuyo fundamento
epistemológico se nutre directamente de la teoría sociológica para
comprender los fenómenos sociales y estructurar modelos de intervención
eficaces. Las corrientes teóricas clásicas y contemporáneas no solo
ofrecen marcos explicativos sobre la organización de la sociedad, sino
que también predeterminan la postura ética, política y metodológica del
profesional en territorio. El presente informe analiza tres de los
paradigmas sociológicos más influyentes: el funcionalismo, el marxismo y
las teorías interpretativas (con énfasis en el interaccionismo
simbólico). El objetivo es desglosar la manera en que cada corriente
conceptualiza la estructura social e influye en el rol, funciones y
alcances del Trabajador Social frente a las problemáticas del entorno
contemporáneo.

\subsection{2. Desarrollo}\label{desarrollo}

\subsubsection{El Funcionalismo y el Rol del Trabajador
Social}\label{el-funcionalismo-y-el-rol-del-trabajador-social}

El funcionalismo, con raíces en el pensamiento de Émile Durkheim,
concibe a la sociedad como un sistema complejo cuyas partes
(instituciones, normas y roles) trabajan de forma interdependiente para
promover la armonía, la cohesión y la estabilidad social
{[}@durkheim2019division{]}. Bajo esta perspectiva, las desviaciones o
conflictos son interpretados como disfunciones temporales del sistema
que requieren mecanismos de adaptación e integración.

Para el Trabajo Social de corte funcionalista, el rol primordial es el
de un facilitador de la adaptación recíproca entre el individuo y su
entorno social. Las funciones del Trabajador Social se centran en el
asistencialismo técnico, la mediación institucional y la rehabilitación
conductual, procurando restablecer el equilibrio sistémico. El
profesional interviene para canalizar recursos del Estado y de las
organizaciones, ayudando a que los sujetos superen sus ``desajustes'' y
se reintegren productivamente a la estructura establecida, sin
cuestionar las bases distributivas u ontológicas de dicha estructura.

\subsubsection{El Marxismo y la Intervención
Crítica}\label{el-marxismo-y-la-intervenciuxf3n-cruxedtica}

A diferencia del consenso funcionalista, el marxismo sitúa al conflicto
material y a la lucha de clases como los motores del desarrollo
histórico y de la dinámica social {[}@marx2018manifiesto{]}. Karl Marx
sostiene que la estructura económica (relaciones de producción)
determina la superestructura política e ideológica, consolidando una
distribución desigual del poder y la riqueza que aliena a las clases
desposeídas.

En este marco, el Trabajo Social Crítico o Liberador rechaza la mera
función asistencialista y adaptativa por considerarla un mecanismo de
control social que reproduce el statu quo. El Trabajador Social marxista
asume un rol político y pedagógico como agente de transformación social.
Sus funciones se orientan a la concienciación de los sujetos sobre sus
condiciones materiales de opresión, la organización comunitaria y la
promoción de la justicia social. La intervención no busca que el
individuo se adapte a un sistema inherentemente desigual, sino
capacitarlo para transformar activamente las estructuras institucionales
y económicas que perpetúan la marginalidad y la vulnerabilidad.

\subsubsection{Las Teorías Interpretativas e Interaccionismo
Simbólico}\label{las-teoruxedas-interpretativas-e-interaccionismo-simbuxf3lico}

Las teorías interpretativas, representadas principalmente por el
interaccionismo simbólico de Herbert Blumer, abandonan las grandes
estructuras macro y se enfocan en las microinteracciones cotidianas y
los procesos de significación subjetiva
{[}@blumer1982interaccionismo{]}. Para esta corriente, la sociedad se
construye y reconstruye constantemente a través de los significados que
los individuos otorgan a sus acciones, símbolos e interacciones en el
día a día.

El Trabajo Social interpretativo y de enfoque construccionista prioriza
el mundo interno de los sujetos {[}@kisnerman1998pensar{]}. El rol del
profesional es el de un copartícipe y facilitador dialógico. Su función
principal consiste en comprender, desde la empatía y la fenomenología,
la realidad de los usuarios basada en sus propias interpretaciones. La
intervención no se impone de manera tecnocrática, sino que se co-diseña
a través del diálogo, ayudando a los sujetos a resignificar sus
trayectorias de vida, reconstruir sus lazos relacionales y desarrollar
procesos de resiliencia y auto-determinación a nivel micro-social.

\subsection{3. Conclusiones}\label{conclusiones}

Los paradigmas analizados demuestran que no existe una única ruta
metodológica en el Trabajo Social, sino una pluralidad epistemológica
que enriquece la práctica. Mientras que el funcionalismo provee
herramientas técnicas para la gestión institucional y la mitigación de
disfunciones inmediatas en pos de la cohesión, el marxismo aporta una
mirada crítica indispensable para desvelar las opresiones estructurales
y guiar la praxis hacia la emancipación colectiva. Por su parte, las
teorías interpretativas devuelven el protagonismo al sujeto, recordando
que la intervención pierde sentido si se ignoran los significados y
vivencias de los actores involucrados. Un Trabajo Social contemporáneo y
ético requiere articular dialécticamente la solidez estructural del
análisis crítico con la sensibilidad relacional del enfoque
interpretativo.

\subsection{4. Bibliografía}\label{bibliografuxeda}

\end{document}
